\documentclass[11pt]{article}

\usepackage[margin=1in]{geometry}
\usepackage[T1]{fontenc}
\usepackage{lmodern}
\usepackage{microtype}
\usepackage{amsmath,amssymb,amsthm,mathtools}
\usepackage{booktabs,tabularx,array,multirow}
\usepackage{graphicx}
\usepackage{caption}
\usepackage{subcaption}
\usepackage{placeins}
\usepackage{xcolor}
\usepackage{enumitem}
\usepackage{natbib}
\usepackage[hidelinks]{hyperref}
\usepackage{fancyhdr}
\usepackage{seqsplit}

\definecolor{navy}{HTML}{17365D}
\definecolor{teal}{HTML}{00897B}
\definecolor{orange}{HTML}{E76F51}
\hypersetup{colorlinks=true,linkcolor=navy,citecolor=teal,urlcolor=navy}
\graphicspath{{figures/}}
\captionsetup{font=small,labelfont=bf}
\setlist{nosep,leftmargin=1.5em}
\setlength{\parskip}{0.35em}
\setlength{\parindent}{1.5em}
\setlength{\emergencystretch}{2em}
\pagestyle{fancy}
\fancyhf{}
\fancyhead[L]{Prefix-Causal Harm-Budget Battery SOH Estimation}
\fancyhead[R]{Yuyang Wu}
\fancyfoot[C]{\thepage}
\renewcommand{\headrulewidth}{0.4pt}
\setlength{\headheight}{14pt}

\newtheorem{lemma}{Lemma}
\newtheorem{theorem}{Theorem}
\newtheorem{corollary}{Corollary}
\newcommand{\R}{\mathbb{R}}

\title{\textbf{Risk-Controlled Cross-Domain Lithium-Ion Battery State-of-Health Estimation via Prefix-Causal Harm-Budget Projection}}
\author{Yuyang Wu}
\date{}

\begin{document}
\maketitle

\begin{abstract}
Battery-management systems in a new fleet can receive partial charging records before reference capacity labels become available. Cross-domain learned updates can then cause negative transfer, while retrospective monotone smoothing can revise already-issued SOH records. We propose prefix-causal harm-budget projection (PCHP), an output operator separating a protected source predictor from a learned candidate. PCHP projects each candidate into a physically bounded, recursively feasible, non-increasing harm tube around a one-sided causal baseline. For every possible outcome, the absolute-loss increase relative to the protected state is at most the budget without distributional assumptions; this intersection is the exact online viability kernel: every admitted output remains extendable, and the projection is the unique closest-candidate solution. Under a fixed contract, the complete causal operator is prefix-wise $\ell_\infty$-nonexpansive: bounded perturbations in either frozen predictor are not amplified. Loss-geometry analysis derives exact metric and bounded squared-loss regions and proves that no nontrivial finite-budget squared-loss update exists over the real line. Across $12$ domains, $586$ cells, and $601{,}932$ records, PCHP reduced domain-equal cell-macro MAE from $0.07524$ to $0.07038$. The paired change was $-0.00486$ with a domain-cluster $95\%$ interval of $[-0.00671,-0.00301]$; $11/12$ domains improved. An unprotected comparator reached $0.05470$ but violated every domain's harm tube, quantifying the accuracy cost of harm control. In outcome-blind, protocol-locked external confirmation, PCHP reduced cell-equal error from $0.17496$ to $0.16510$ across $45$ cells; all improved. A NASA stress test retained the deterministic certificate but exposed target incompatibility. PCHP therefore enables auditable, risk-bounded online battery-health adaptation before reference capacity becomes available.
\end{abstract}

\noindent\textbf{Keywords:} lithium-ion battery; state of health; cross-domain learning; prefix causality; deterministic harm bound; monotone prediction; external validation

\section{Introduction}

Consider a battery-swapping operator that introduces cells from a new supplier. During the first months of service, the battery-management system receives partial charging records, while an accurate capacity label remains unavailable because the operator has not interrupted service for a complete reference characterization. This information mismatch is operationally realistic: conventional capacity estimators often require specially designed charge--discharge conditions, whereas recent methods use random or partial charging segments because complete charging profiles seldom match field operation \citep{liu2023rapid,zhang2024flexible}. A source-trained model can estimate SOH immediately. A second model can exploit how the cell's charge response has changed since commissioning. Yet neither the direction nor the magnitude of that update can be validated on the new fleet before outcomes arrive. Worse, if non-increasing SOH is imposed retrospectively, a downward anomaly observed next month can revise the estimate that was reported today. Such a system is unsuitable for an online maintenance record even when its final reconstructed trajectory is smooth.

SOH estimation is central to maintenance and warranty decisions, and data-driven models can exploit voltage, current, charge, temperature, and cycle history without identifying every degradation pathway \citep{ng2020mlbattery,roman2021pipeline}. Short charge segments can support SOH estimation, while early-life records can support lifetime prediction before full lifetime characterization \citep{roman2021pipeline,severson2019cyclelife}. Cross-domain learning further reduces reliance on labelled ageing experiments through representation alignment, source selection, personalized transfer, ensembles, and test-time adaptation \citep{li2020sstca,ma2022personalized,lu2023without,duan2025wasserstein,feng2024batteryttt}. These advances address how to generate a useful update. They do not, by themselves, answer two deployment questions: whether the update can be bounded relative to a protected predictor for every hidden outcome, and whether a trajectory constraint can be executed without future access.

We address this capability gap by separating \emph{utility generation}, \emph{causal state construction}, and \emph{harm control}. A protected model uses raw early-charge features; a learned candidate additionally uses raw-scale changes from the first five characterization records. The protected raw predictions are converted into a prefix-causal non-increasing state that ignores upward innovations and incorporates a fraction $\alpha$ of downward innovations. At each record, PCHP projects the current candidate into a recursively feasible interval around that state. The construction requires no target outcome, never revises an emitted prefix, and limits the worst possible increase in absolute loss relative to the protected state. Figure~\ref{fig:workflow} summarizes the operational setting and evidence chain.

\begin{figure}[t]
  \centering
  \includegraphics[width=\textwidth]{fig1_method_workflow.pdf}
  \caption{Operational problem, prefix-causal method, and headline evidence. \textbf{A}, a deployed SOH record must not change when future charge measurements arrive. \textbf{B}, a protected raw prediction is converted into a one-sided causal non-increasing state, while a commissioning-referenced model proposes a candidate update. \textbf{C}, PCHP projects the candidate into the intersection of the harm tube and the recursively feasible non-increasing set. \textbf{D}, paired effects from nested development, the candidate-free adaptive-value control, and outcome-blind prospectively locked external validation; intervals resample the stated independent units.}
  \label{fig:workflow}
\end{figure}

The single core contribution is a recursively feasible projection that joins three properties usually treated separately: prefix-only execution, non-increasing output, and a pointwise harm budget around a protected prediction. Its theory characterizes the harm tube as the necessary-and-sufficient current loss region and its trajectory-constrained intersection as the exact online viability kernel: every admitted output remains extendable under every future non-increasing protected-state path, whereas every excluded output already violates the contract. The complete two-predictor causal cascade is additionally prefix-wise $\ell_\infty$-nonexpansive with tight constant one, so its recursion does not magnify bounded perturbations in the protected or candidate stream. The analysis also proves that unconditional zero-harm updating is impossible unless the update leaves the protected prediction unchanged and identifies PCHP as the unique closest-candidate solution in that kernel. It further exposes a loss-dependent feasibility boundary: metric losses induce exact metric balls; squared loss on a bounded SOH range induces an exact baseline-dependent asymmetric interval; and squared loss over the full real line makes every nontrivial finite-budget update impossible. A directionally weighted corollary converts unequal under- and overestimation costs into an exact asymmetric tube; for time-varying costs and budgets, the same analysis gives an exact recursive-feasibility condition and a tight minimum current budget. These results remain one output-contract contribution. The commissioning-referenced candidate is its utility mechanism, not a second architectural contribution. We support the method through five linked evidence layers: a future-extension falsification showing that PCHP never revises the past; nested twelve-domain evaluation with source-only $\alpha$ selection and a candidate-free harm-budget-shift control; a source-tuned unprotected comparator that exposes the utility sacrificed by the declared budget; an outcome-blind, prospectively locked Bayreuth external validation with predictions fixed before capacity access; and a protocol-frozen NASA stress test that retains its scientific failure boundary when discharge protocols make the absolute target incomparable.

The guarantee is intentionally relative. PCHP does not certify that the protected baseline is accurate, that the candidate improves every domain, that monotone SOH captures short-term capacity recovery, or that $\delta=0.01$ is appropriate for every operational decision. It certifies that an emitted update is prefix causal, non-increasing, and cannot increase absolute loss by more than the declared budget relative to its protected causal state.

\section{Related work and novelty boundary}

\subsection{Cross-domain battery health estimation}

Battery SOH estimation has progressed from handcrafted health indicators and shallow regressors toward sequence models, physics-informed learning, and cross-domain transfer \citep{ng2020mlbattery,roman2021pipeline,wang2024pinn}. Semi-supervised transfer component analysis aligns source and target representations with limited labelled data \citep{li2020sstca}; deep transfer learning personalizes health prediction under unseen discharge protocols \citep{ma2022personalized}; and domain-adaptive ensembles estimate SOH without additional target degradation experiments \citep{lu2023without}. Recent work chooses a source domain using Wasserstein distance before fine-tuning a temporal model \citep{duan2025wasserstein}, while BatteryTTT adapts continually using arriving unlabelled target inputs \citep{feng2024batteryttt}. In Applied Energy, multi-source adaptation has combined source selection with ageing-order-aware alignment \citep{qiu2024multisource}, while Bayesian mixed-effects transfer adapts parameter updating to the operating information available for each battery \citep{zhang2025bayesian}. A complementary Energy and AI field study combines multi-view co-training with confidence-driven pseudo-labeling across $3{,}000$ vehicles and incorporates labels incrementally as they become available \citep{hadzalic2025field}. These studies show how target history, probabilistic information sharing, or progressive labels can improve estimation under heterogeneity. Our question begins before any target outcome is available: can a learned update be permitted while bounding its relative harm for every hidden outcome and preventing future revision of issued records?

Concurrent methods extend the adaptation direction through domain-adversarial cell-to-pack transfer with limited target supervision \citep{shu2026celltopack}, probabilistic transfer under sparse target labels \citep{shao2026probabilistic}, and physics-guided expert routing with a global monotonicity penalty \citep{jiang2026phasemoe}. Reliability-oriented extensions combine domain alignment with conformal prediction intervals \citep{filgueira2026conformalized} or select a monotone post-hoc calibrator through labelled field holdout error \citep{nyachionjeka2026safecalibrated}. These methods improve adaptation, uncertainty quantification, or physical consistency, but they do not supply the guarantee studied here. Our target domain contributes neither outcomes nor parameter updates: two source-only predictors are frozen, only their output trajectories are constrained, and the guarantee concerns a pointwise loss increase rather than interval coverage, target-supervised improvement, or a training penalty.

Adjacent prognostics and health-management studies sharpen this distinction. Compound-domain-shift prognostics aligns both inter-domain distributions and degradation-period distributions within machinery trajectories \citep{ding2023compound}. Source-only open-set domain generalization learns fault classifiers for unknown working conditions when target-domain data are absent during training \citep{song2026regroupment}. Interactive battery prognostics recursively couples neural capacity estimates with a switching state-space process to update remaining useful life \citep{duan2025interactive}. Their reported objectives are, respectively, representation-shift reduction, source-only open-set classification, and recursive life prediction. PCHP targets a different object: an outcome-agnostic, pointwise harm budget around a protected prefix-causal state while learned predictors remain frozen.

The BatteryML and BatteryLife projects standardize heterogeneous public ageing data and make broad cross-dataset evaluation more reproducible \citep{zhang2024batteryml,tan2025batterylife}. Such breadth creates a statistical responsibility: cycles are nested within cells and cells are nested within datasets. We therefore treat a complete battery dataset domain, not a cycle, as the independent unit for the development cross-domain effect.

\subsection{Negative transfer and safe regression}

Negative transfer includes cases in which transferred data or knowledge reduce target performance, and proposed remedies include domain-similarity estimation, safe transfer, and explicit mitigation \citep{zhang2023negative}. In safe semi-supervised regression, SAFER geometrically projects a supervised prediction onto a convex set of semi-supervised regressors and obtains a no-worse squared-loss result when the complete truth vector lies in that set \citep{li2017safer}. SAFEW extends maximin safe weak supervision to several convex classification and regression losses, again under a condition linking the complete truth vector to the base-learner hull \citep{li2021safeweak}. Prediction with expert evaluators permits proper mixable losses that vary across experts and time, but guarantees cumulative regret after outcomes arrive \citep{chernov2009evaluators}; online probabilistic recalibration likewise pursues calibration together with vanishing or low cumulative regret relative to baseline forecasters \citep{marx2025calibrated,deshpande2025calibrated}. Conformal calibrators establish calibration in probability and analyze an efficiency form of ``do no harm'' in which conformalizing an ideal predictive system should not make its predictive distribution substantially worse \citep{vovk2020conformal}. Weighted conformal transport instead uses unlabelled target covariates to recover marginal coverage under a covariate-shift model \citep{tuwani2023transport}, and battery-specific conformal transfer combines representation alignment with calibrated prediction intervals \citep{filgueira2026conformalized}. Protected Online Entropy Matching detects classification-entropy shifts and adapts model parameters so that accuracy and calibration are retained empirically when no shift is present \citep{bar2024protected}. Labelled battery field data have also been used to choose a monotone post-hoc calibrator only when held-out error improves \citep{nyachionjeka2026safecalibrated}. These are substantive protection mechanisms, but their protected objects are conditional vector risk, cumulative regret, predictive-distribution efficiency, marginal coverage, empirical holdout error, or an adaptive training process---not the loss of each emitted point under every hidden outcome.

Our result differs in guarantee object and execution contract. We do not claim that the final prediction is always no worse. Instead, for every real-valued truth and without a truth-in-set condition, we place an exact upper bound on the largest possible increase in absolute loss at each emitted record. The guarantee is checked before that truth is observed, both predictors remain frozen, and the same projection must preserve every previously emitted non-increasing prefix. Table~\ref{tab:novelty_audit} compares these information and guarantee boundaries directly.

\begin{table}[p]
  \centering
  \caption{Closest-method audit by guarantee object and information boundary. Conditional records the principal condition under which the cited protection is stated; it is not a judgment of method quality.}
  \label{tab:novelty_audit}
  \footnotesize
  \renewcommand{\arraystretch}{0.94}
  \begin{tabularx}{\textwidth}{@{}>{\raggedright\arraybackslash}p{0.17\textwidth}>{\raggedright\arraybackslash}p{0.20\textwidth}>{\raggedright\arraybackslash}p{0.25\textwidth}>{\raggedright\arraybackslash}X@{}}
    \toprule
    Approach & Protected object & Information or condition & Distinction from PCHP \\
    \midrule
    SAFER / SAFEW \citep{li2017safer,li2021safeweak} & Full-vector loss relative to a supervised baseline & Conditional on the truth vector lying in the base-regressor hull & Offline vector or maximin prediction; no per-record outcome-uniform budget or trajectory recursion \\
    Expert advice / online recalibration / conservative OCO \citep{chernov2009evaluators,marx2025calibrated,deshpande2025calibrated,bernasconi2021conservative} & Calibration, cumulative regret, or cumulative default-relative loss & Outcomes arrive sequentially; conservative OCO spends realized loss budget & Cumulative guarantee rather than a same-record outcome-uniform harm region \\
    Conformal calibration / transport \citep{vovk2020conformal,tuwani2023transport,filgueira2026conformalized} & Calibration, ideal-system efficiency, or marginal set coverage & IID ideal-base analysis or the stated covariate-shift assumptions & Protects predictive distributions or sets; no arbitrary-baseline point-loss budget or immutable trajectory \\
    Protected adaptation \citep{bar2024protected} & Accuracy and calibration in no-shift streams & Target entropy stream updates model parameters & Empirical retention around a training process, not a frozen-model deterministic pointwise certificate \\
    Battery safe calibration \citep{nyachionjeka2026safecalibrated} & Held-out target error and monotone scaling & Labelled target holdout selects the calibrator & Target-supervised post-hoc mapping; no outcome-uniform pre-label certificate \\
    Baseline-relative / constrained online decisions \citep{moradipari2020stagewise,deb2025conservative,sridharan2025unknown,vaskov2024donoharm} & Cumulative policy loss, instantaneous expected reward, or action feasibility & Action or constraint feedback, exploration, learned constraints, or counterfactual evaluation & Protects decisions under policy or constraint models, not every-outcome point-prediction loss or issued prefixes \\
    Predictive safety filters \citep{wabersich2021predictive} & Dynamical state and input constraints & Plant model, uncertainty set, backup policy, and terminal safe set & Minimally modifies control actions; no hidden-outcome relative-loss region or immutable prediction record \\
    Monotonic-walk / moving-horizon filters \citep{samar2004moving,gorinevsky2008efficient} & Latent monotone deterioration trend and filter stability & One-sided process model or moving-window constrained optimization & No frozen two-predictor outcome-uniform harm budget or immutable emitted prefix \\
    PCHP (this work) & Per-record absolute-loss increase from a protected causal state & No target labels, outcome distribution, or independence assumption & Exact budgeted tube joined to prefix causality and recursive non-increasing feasibility \\
    \bottomrule
  \end{tabularx}
\end{table}

The closest conceptual parallels in Table~\ref{tab:novelty_audit} are conservative online optimization and baseline-relative sequential decisions. Conservative Projection maps an online convex optimizer's proposal into a default-centred ball so that cumulative realized loss stays within a multiplicative allowance at every round \citep{bernasconi2021conservative}. PCHP instead certifies the current prediction against every possible value of its still-hidden outcome; it cannot borrow surplus from earlier records. Stage-wise conservative linear bandits impose an instantaneous expected-reward floor relative to a baseline action \citep{moradipari2020stagewise}, while general online learning with unknown constraints maps oracle proposals to actions satisfying a learned constraint with high probability \citep{sridharan2025unknown}. Conservative contextual bandits protect cumulative policy loss \citep{deb2025conservative}, and counterfactual do-no-harm control bounds policy-caused constraint violations relative to an alternate safe policy \citep{vaskov2024donoharm}. These are important per-round or anytime protections, but they concern actions, expected rewards, learned constraints, or cumulative feedback. They do not provide an outcome-uniform relative-loss region for an emitted point prediction, nor prevent later records from revising an issued prediction prefix.

Predictive safety filters provide a closer algorithmic analogy: they minimally modify a proposed learning-based control input so that a model-predictive backup policy preserves state and input constraints \citep{wabersich2021predictive}. PCHP likewise filters a learned proposal, but its admissible object is an emitted prediction rather than a control action, its uncertainty is the hidden outcome rather than plant disturbance, and its viability set is available in closed form from relative loss and record-order constraints without a plant model, rollout, or terminal safe set.

The novelty is therefore not scalar clipping, geometric projection, or the discovery of negative transfer. The audited distinction is the battery-deployment conjunction of a learned source-only candidate, a prefix-causal protected state, a recursively feasible monotone output, and an exact pointwise absolute-loss harm tube. Within the literature examined, we did not identify this conjunction; this scoped distinction does not establish universal priority.

The loss-geometry extension is bounded in the same way. The metric identity follows from established metric properties and is not claimed as a new mathematical primitive. The narrower distinction is the exact outcome-range feasibility boundary for squared loss and its integration into the same prefix-causal recursive contract; safe weak-supervision methods condition vector-level guarantees on a truth set, while expert-advice and recalibration methods protect cumulative loss after outcomes arrive.

Recursive feasibility is a longstanding constrained-control concept; its strong form requires feasibility under every feasible input sequence \citep{lofberg2012oops}. Time-varying and history-coupled constraints have also been studied in predictive control and online optimization \citep{li2021information,garber2024projectionfree,knaup2026recursively}. Those works concern dynamical states, cumulative constraint violation, dynamic regret, or chance constraints. They do not establish the outcome-uniform pointwise loss region studied here. We therefore make no priority claim for recursive feasibility or time-varying constraints; our narrower result is the exact condition induced by the scalar PCHP harm-tube recursion.

\subsection{Monotone degradation trends}

Long-term capacity fade motivates non-increasing trend constraints, but measured capacity can exhibit local recovery after rest and other fluctuations \citep{ma2021regeneration}. Moving-horizon monotonic-trend filters formulate a constrained quadratic program over a finite active window \citep{samar2004moving}; monotonic-walk filtering further develops a probabilistic moving-horizon formulation and analyzes bounded-input bounded-output stability and nonlinear impulse-response decay \citep{gorinevsky2008efficient}. These are close stability precedents, but an active window may be recomputed as new observations arrive, and their protected object is a latent statistical trend. Online isotonic regression instead predicts sequentially and is analyzed against the best isotonic function in hindsight, including extensions beyond squared loss \citep{kotlowski2016online}; its objective is adversarial online regret rather than protection around a battery-specific baseline. PCHP therefore does not claim the first stable monotone filter. Its narrower result is incremental prefix-level nonexpansiveness of a frozen two-predictor, moving-harm-tube cascade whose emitted prefix is immutable. Battery studies also use the term ``online SOH estimation'' for inference from a short charge or discharge segment \citep{liu2020online}, but online feature availability does not by itself prevent a later trajectory smoother from revising past outputs. PCHP addresses that operational gap. Its monotone state estimates a latent long-term trend and may suppress real short-term capacity recovery, so raw measurements and the protected trend should remain separately auditable.

\section{Problem formulation}

Let $i$ index a physical cell and $t$ its ordered characterization records. The capacity-based target is
\begin{equation}
  y_{it}=\frac{Q_{it}}{\operatorname{median}(Q_{i1},\ldots,Q_{i5})},
  \label{eq:soh}
\end{equation}
where the first five valid characterization capacities define a cell-specific outcome reference. During prediction, neither $Q_{it}$ nor $y_{it}$ is available. Predictions begin after the first five charge records, whose measurements---but not their discharge capacities---also define the commissioning feature reference. A dataset domain groups cells from one experimental source, instrument family, chemistry, or protocol. In each development fold, the complete target domain is excluded from fitting and hyperparameter selection.

Let $\mathcal F_{it}$ contain the frozen source models, the first-five charge reference, and charge measurements observed for cell $i$ through record $t$. An online prediction is prefix causal if it is measurable with respect to $\mathcal F_{it}$ and is unchanged by any extension after $t$. Let $r_{it}$ be the protected model's raw prediction and $c_{it}$ the learned candidate, both $\mathcal F_{it}$-measurable. The desired output $p_{it}$ must satisfy four requirements: prefix causality, a non-increasing trajectory, a physical range $\mathcal Y=[0,1.3]$, and an absolute-loss harm budget $\delta$ relative to a protected causal state $b_{it}$.

The protected state is a contract, not an assertion of optimality. Its accuracy determines the neighborhood in which PCHP can operate; the theorem cannot turn a poor baseline or an incompatible target into an accurate estimator.

\section{Method}

\subsection{Early-charge and commissioning-referenced features}

Each record is summarized from the first $600\,\mathrm{s}$ of charge. The $24$ raw channels include normalized accumulated charge, voltage mean and standard deviation, voltage slope, current statistics, temperature availability and statistics, voltages at $60\,\mathrm{s}$ intervals, and currents at $0$, $300$, and $600\,\mathrm{s}$. The log-transformed cycle number is appended. Missing values are median-imputed from the source training fold, with missingness indicators added by the imputer.

For raw channel $j$, the commissioning reference for cell $i$ is
\begin{equation}
  a_{ij}=\operatorname{median}_{1\leq s\leq5}x_{ijs},
\end{equation}
and the signed raw-scale change feature is
\begin{equation}
  d_{ijt}=x_{ijt}-a_{ij}.
  \label{eq:change}
\end{equation}
The protected model receives the raw features and cycle number. The candidate receives the same inputs plus all $d_{ijt}$. Normalized relative changes were considered during development but removed before external access because they increased domain-equal error and amplified small-denominator noise. This removal is a development ablation, not an externally confirmed discovery.

Both primary predictors are ExtraTrees regressors with $80$ trees, minimum leaf size $20$, feature subsampling $0.8$, and fixed seed $20{,}260{,}801$. Source sampling is capped at $80$ temporally spaced records per physical cell. The baseline uses ordinary pooled weights; the candidate assigns each source domain equal total weight so that large archives cannot dominate merely by contributing more cycles. No target-domain outcome, fine-tuning, or source selection occurs at inference.

A protocol-locked secondary audit replaced both regressors with Ridge, HistGradientBoosting, or LightGBM while preserving the same data, features, weighting, outer-domain exclusions, harm budget, PCHP operator, and exact V326 outer-domain $\alpha$ schedule. No secondary-backbone hyperparameter or $\alpha$ tuning was allowed. Because all development outcomes had already been opened, this audit tests zero-additional-tuning portability rather than external confirmation or model selection.

\subsection{Prefix-causal protected state}

Let $\Pi_{\mathcal Y}$ denote scalar clipping to $\mathcal Y$. For one physical cell, PCHP initializes and updates the protected state as
\begin{align}
  b_{i1}&=\Pi_{\mathcal Y}(r_{i1}),\\
  b_{it}&=\Pi_{\mathcal Y}\!\left[b_{i,t-1}
  +\alpha\min\{r_{it}-b_{i,t-1},0\}\right],\qquad t>1,
  \label{eq:causal_state}
\end{align}
where $\alpha\in(0,1]$ is selected using source domains only. Upward innovations are ignored; a downward innovation is assimilated at rate $\alpha$. Thus $b_{it}\leq b_{i,t-1}$, while an isolated negative shock is attenuated when $\alpha<1$.

\begin{lemma}[Prefix invariance of the protected state]
\label{lem:prefix}
For any $k\leq T$, computing Eq.~\eqref{eq:causal_state} on $(r_{i1},\ldots,r_{ik})$ gives the same first $k$ states as computing it on any extension $(r_{i1},\ldots,r_{iT})$.
\end{lemma}
\begin{proof}
The initialization uses only $r_{i1}$, and each subsequent state is a deterministic function of $(b_{i,t-1},r_{it})$. Induction over $t$ proves both prefix invariance and $\mathcal F_{it}$-measurability.
\end{proof}

For $t>1$, the same update has the local proximal form
\begin{equation}
  b_{it}=\underset{q\in\mathcal Y,\ q\leq b_{i,t-1}}{\arg\min}
  \left\{
  \frac{\alpha}{2}(q-r_{it})^2+
  \frac{1-\alpha}{2}(q-b_{i,t-1})^2
  \right\}.
  \label{eq:protected_proximal}
\end{equation}
Thus $\alpha$ controls current-prediction fidelity against state inertia under the one-sided order constraint. This is an exact algorithmic characterization, not a claim of statistical or electrochemical optimality.

\subsection{Recursively feasible harm-budget projection}

Set $p_{i0}=1.3$. At record $t$, define
\begin{equation}
  I_{it}=\left[
  \max\{0,b_{it}-\delta\},
  \min\{1.3,b_{it}+\delta,p_{i,t-1}\}
  \right]
  \label{eq:feasible_interval}
\end{equation}
The output is
\begin{equation}
  p_{it}=\Pi_{I_{it}}(c_{it}).
  \label{eq:method}
\end{equation}
The upper endpoint enforces the previously emitted prediction, whereas the tube around $b_{it}$ limits displacement from the protected state. The projection is the closest currently feasible value to the learned candidate.

\begin{theorem}[Recursive feasibility, prefix causality, and monotonicity]
\label{thm:pchp}
If $b_{it}$ follows Eq.~\eqref{eq:causal_state} and $\delta\geq0$ is constant, then $I_{it}$ is nonempty for every $t$. The output in Eq.~\eqref{eq:method} is prefix causal, lies in $\mathcal Y$, is non-increasing in $t$, and satisfies $|p_{it}-b_{it}|\leq\delta$.
\end{theorem}
\begin{proof}
At $t=1$, $p_{i0}=1.3$ and $b_{i1}\in\mathcal Y$, so the interval is nonempty. Suppose it was feasible at $t-1$. Because $b_{it}\leq b_{i,t-1}$ and $p_{i,t-1}\geq\max\{0,b_{i,t-1}-\delta\}$, we have $p_{i,t-1}\geq\max\{0,b_{it}-\delta\}$; hence the new lower endpoint does not exceed any term defining the upper endpoint. Projection into $I_{it}$ gives the range and displacement properties, and $p_{it}\leq p_{i,t-1}$ gives monotonicity. Every quantity uses only the observed prefix, so the output is prefix causal.
\end{proof}

\begin{theorem}[Prefix-level nonexpansiveness of the complete operator]
\label{thm:nonexpansive}
Consider two runs for the same cell with common $\alpha$, $\delta$, $\mathcal Y$, and $p_{i0}=\widetilde p_{i0}=1.3$, but protected-prediction streams $r,\widetilde r$ and candidate streams $c,\widetilde c$. Let $b,\widetilde b$ and $p,\widetilde p$ be their corresponding protected states and PCHP outputs. For every prefix ending at $t$,
\begin{align}
  \max_{1\leq s\leq t}|b_{is}-\widetilde b_{is}|
  &\leq \max_{1\leq s\leq t}|r_{is}-\widetilde r_{is}|,\\
  \max_{1\leq s\leq t}|p_{is}-\widetilde p_{is}|
  &\leq
  \max\!\left\{
  \max_{1\leq s\leq t}|r_{is}-\widetilde r_{is}|,
  \max_{1\leq s\leq t}|c_{is}-\widetilde c_{is}|
  \right\}.
  \label{eq:prefix_nonexpansive}
\end{align}
The Lipschitz constant one is tight.
\end{theorem}
\begin{proof}
Write
\[
F_\alpha(x,z)=\Pi_{\mathcal Y}\!\left((1-\alpha)x+\alpha\min\{x,z\}\right).
\]
The minimum is jointly one-Lipschitz in the sup norm; convex combination with $x$ and scalar clipping preserve that constant. Initialization and induction therefore give the first inequality. For the projection step, let
\[
L(b)=\max\{0,b-\delta\},\qquad
U(b,u)=\min\{1.3,b+\delta,u\}.
\]
Equation~\eqref{eq:method} is equivalently $p_{it}=\operatorname{med}\{L(b_{it}),c_{it},U(b_{it},p_{i,t-1})\}$. Minimum, maximum, and median are jointly one-Lipschitz in the sup norm, so
\[
|p_{it}-\widetilde p_{it}|
\leq \max\{|b_{it}-\widetilde b_{it}|,
|c_{it}-\widetilde c_{it}|,
|p_{i,t-1}-\widetilde p_{i,t-1}|\}.
\]
Induction and the protected-state bound prove Eq.~\eqref{eq:prefix_nonexpansive}. A one-step configuration with identical protected predictions and two interior candidates separated by $\varepsilon>0$ produces outputs separated by exactly $\varepsilon$, establishing tightness.
\end{proof}

\subsection{Exact loss-dependent harm geometry}

For a prediction space $\mathcal P$, outcome set $\mathcal Z$, loss $\ell$, protected prediction $b$, and budget $\delta\geq0$, define the outcome-uniform harm region
\begin{equation}
  \mathcal H^{\ell}_{\delta}(b;\mathcal Z)
  =\left\{p\in\mathcal P:
  \sup_{y\in\mathcal Z}\bigl[\ell(p,y)-\ell(b,y)\bigr]\leq\delta
  \right\}.
  \label{eq:generic_harm_region}
\end{equation}
The outcome set is part of the contract: changing it can change whether any nontrivial update is certifiable.

\begin{theorem}[Exact metric-loss geometry]
\label{thm:metric_harm}
Let $(\mathcal M,d)$ be any metric space and set $\ell(q,y)=d(q,y)$ for $q,y\in\mathcal M$. For arbitrary $b,p\in\mathcal M$,
\begin{equation}
  \sup_{y\in\mathcal M}\left[d(p,y)-d(b,y)\right]=d(p,b),
  \label{eq:metric_identity}
\end{equation}
and therefore
\begin{equation}
  \mathcal H^{d}_{\delta}(b;\mathcal M)
  =\{p\in\mathcal M:d(p,b)\leq\delta\}.
  \label{eq:metric_ball}
\end{equation}
\end{theorem}
\begin{proof}
The triangle inequality gives $d(p,y)-d(b,y)\leq d(p,b)$ for every $y$. Equality is attained at $y=b$. Solving the resulting inequality gives Eq.~\eqref{eq:metric_ball}.
\end{proof}

Theorem~\ref{thm:metric_harm} is a metric-space consequence of the triangle inequality, not a priority claim for that mathematical primitive. Its role is to show that exact outcome-uniform harm is a geometry, not an absolute-loss coincidence. Scalar absolute loss is the special case $\mathcal M=\R$ and $d(p,y)=|p-y|$.

\begin{theorem}[Exact harm region and zero-budget impossibility]
\label{thm:harm_region}
For fixed $b\in\R$ and $\delta\geq0$, define
\begin{equation}
  \mathcal H_\delta(b)=\left\{p\in\R:
  \sup_{y\in\R}\bigl(|p-y|-|b-y|\bigr)\leq\delta\right\}.
  \label{eq:harm_region}
\end{equation}
Then $\mathcal H_\delta(b)=[b-\delta,b+\delta]$. Consequently,
\begin{equation}
  I_{it}=\mathcal H_\delta(b_{it})\cap\mathcal Y
  \cap(-\infty,p_{i,t-1}],
  \label{eq:maximal_interval}
\end{equation}
so $I_{it}$ is exactly the largest current-record output set that simultaneously satisfies the physical range, the previously emitted non-increase constraint, and the declared worst-case absolute-loss budget. In particular, $\mathcal H_0(b)=\{b\}$: any nontrivial outcome-free update $p\neq b$ worsens absolute loss for some possible outcome by exactly $|p-b|$.
\end{theorem}
\begin{proof}
Theorem~\ref{thm:metric_harm} makes the defining inequality in Eq.~\eqref{eq:harm_region} equivalent to $|p-b|\leq\delta$, proving the interval identity. Intersecting this exact harm region with $\mathcal Y=[0,1.3]$ and $p\leq p_{i,t-1}$ gives Eq.~\eqref{eq:maximal_interval}. Setting $\delta=0$ leaves only $p=b$; for $p\neq b$, Eq.~\eqref{eq:metric_identity} supplies an outcome attaining harm $|p-b|>0$.
\end{proof}

\begin{corollary}[Exact online viability kernel]
\label{cor:viability_kernel}
Fix a current prefix, a constant budget $\delta\geq0$, and a non-increasing protected state. Call a current output $q$ robustly extendable if it satisfies the current range, non-increase, and outcome-uniform harm constraints and, for every finite future non-increasing protected-state continuation in $\mathcal Y$, admits an output continuation satisfying the same constraints. The set of all robustly extendable current outputs is exactly $I_{it}$ in Eq.~\eqref{eq:feasible_interval}. Hence $I_{it}$ is the exact online viability kernel of the declared output contract, and PCHP is the unique closest-candidate point in that kernel.
\end{corollary}
\begin{proof}
Every $q\in I_{it}$ satisfies the current constraints by Theorem~\ref{thm:harm_region}. For any future non-increasing continuation $b_{is}$, $s>t$, define $L_{is}=\max\{0,b_{is}-\delta\}$. Then $L_{is}$ is non-increasing, lies in $\mathcal Y$, satisfies $|L_{is}-b_{is}|\leq\delta$, and obeys $L_{i,t+1}\leq L_{it}\leq q$. Thus $p_{is}=L_{is}$ is a feasible continuation for every such protected-state path. Conversely, any $q\notin I_{it}$ already violates the current range, issued-prefix order, or exact harm constraint and therefore cannot be robustly extendable. Corollary~\ref{cor:optimality} gives the unique closest-candidate statement.
\end{proof}

Thus the tube is not a heuristic sufficient condition: it is the exact feasible region for the stated pointwise claim. Combining Theorem~\ref{thm:harm_region} with Theorem~\ref{thm:pchp} yields, record by record,
\begin{equation}
  |p_{it}-y_{it}|-|b_{it}-y_{it}|
  \leq |p_{it}-b_{it}|\leq\delta.
  \label{eq:pointwise_harm}
\end{equation}
The recursive monotonicity constraint can make $I_{it}$ narrower than the standalone harm tube.

\begin{theorem}[Squared-loss outcome-range dichotomy]
\label{thm:squared_geometry}
Let $L<U$, let $b,p\in[L,U]$, and define $\ell_2(q,y)=(q-y)^2$. If the declared outcome set is $[L,U]$, then
\begin{equation}
  \sup_{y\in[L,U]}\left[(p-y)^2-(b-y)^2\right]
  =
  \begin{cases}
    (p-b)(p+b-2L), & p\geq b,\\
    (b-p)(2U-p-b), & p<b.
  \end{cases}
  \label{eq:squared_identity}
\end{equation}
Consequently, the exact finite-budget region is
\begin{equation}
  \mathcal H^{(2)}_{\delta}(b;[L,U])
  =\left[
  \max\!\left\{L,U-\sqrt{(U-b)^2+\delta}\right\},
  \min\!\left\{U,L+\sqrt{(b-L)^2+\delta}\right\}
  \right].
  \label{eq:squared_region}
\end{equation}
If the declared outcome set is instead $\R$, then for every $p\neq b$,
\begin{equation}
  \sup_{y\in\R}\left[(p-y)^2-(b-y)^2\right]=+\infty.
  \label{eq:squared_unbounded}
\end{equation}
Thus, when the outcome set is the full real line, no nontrivial update has a finite outcome-uniform squared-loss budget.
\end{theorem}
\begin{proof}
Expanding the loss difference gives $(p-b)(p+b-2y)$. When $p\geq b$, this affine function is maximized at $y=L$; when $p<b$, it is maximized at $y=U$, proving Eq.~\eqref{eq:squared_identity}. On the upper side of $b$, the budget inequality is $(p-L)^2-(b-L)^2\leq\delta$; on the lower side it is $(U-p)^2-(U-b)^2\leq\delta$. Solving these inequalities and intersecting with $[L,U]$ gives Eq.~\eqref{eq:squared_region}. If $p>b$, the expanded difference diverges as $y\to-\infty$; if $p<b$, it diverges as $y\to+\infty$, proving Eq.~\eqref{eq:squared_unbounded}.
\end{proof}

\begin{corollary}[Recursive feasibility under bounded squared loss]
\label{cor:squared_recursive}
For fixed $\delta\geq0$, define
\begin{align}
  \underline q^{(2)}_{\delta}(b)
  &=\max\!\left\{L,U-\sqrt{(U-b)^2+\delta}\right\},\\
  \overline q^{(2)}_{\delta}(b)
  &=\min\!\left\{U,L+\sqrt{(b-L)^2+\delta}\right\}.
  \label{eq:squared_bounds}
\end{align}
If $b_{it}\in[L,U]$ is prefix causal and non-increasing, set $p_{i0}=U$ and
\begin{equation}
  I^{(2)}_{it}
  =\left[
  \underline q^{(2)}_{\delta}(b_{it}),
  \min\!\left\{\overline q^{(2)}_{\delta}(b_{it}),p_{i,t-1}\right\}
  \right].
  \label{eq:squared_interval}
\end{equation}
Then $I^{(2)}_{it}$ is nonempty for every record and every candidate sequence. Projecting $c_{it}$ into this interval is prefix causal, remains in $[L,U]$, is non-increasing, and satisfies
\begin{equation}
  (p_{it}-y)^2-(b_{it}-y)^2\leq\delta
  \qquad\text{for every }y\in[L,U].
  \label{eq:squared_pointwise_harm}
\end{equation}
\end{corollary}
\begin{proof}
The lower endpoint $\underline q^{(2)}_{\delta}(b)$ is nondecreasing in $b$. Hence $b_{it}\leq b_{i,t-1}$ implies $\underline q^{(2)}_{\delta}(b_{it})\leq\underline q^{(2)}_{\delta}(b_{i,t-1})\leq p_{i,t-1}$. The standalone region in Eq.~\eqref{eq:squared_region} contains $b_{it}$, so its lower endpoint also does not exceed its upper endpoint. The recursive interval is therefore nonempty. Projection gives the range and non-increase properties, prefix measurability gives causality, and Theorem~\ref{thm:squared_geometry} gives Eq.~\eqref{eq:squared_pointwise_harm}.
\end{proof}

This dichotomy explains the experimental loss choice rather than claiming that squared loss is empirically inferior. Absolute loss gives a finite exact certificate without bounding the hidden truth. The bounded-interval squared-loss construction is valid when $[L,U]$ is a credible outcome contract; Theorem~\ref{thm:squared_geometry} separately establishes impossibility over $\R$. Deterministic symbolic, boundary, exterior-point, and random-trajectory falsification tests of Eqs.~\eqref{eq:squared_identity}--\eqref{eq:squared_interval} are reported in the supplement.

\begin{corollary}[Exact asymmetric decision-loss extension]
\label{cor:asymmetric}
Let $c_{\mathrm{under}},c_{\mathrm{over}}>0$ and define the directionally weighted absolute loss
\begin{equation}
  \ell_{\mathrm{asym}}(p,y)
  =c_{\mathrm{under}}(y-p)_{+}
  +c_{\mathrm{over}}(p-y)_{+},
  \label{eq:asymmetric_loss}
\end{equation}
where $(z)_{+}=\max\{z,0\}$. For any $b,p\in\R$,
\begin{equation}
  \sup_{y\in\R}\!\left[
  \ell_{\mathrm{asym}}(p,y)-\ell_{\mathrm{asym}}(b,y)
  \right]
  =
  \begin{cases}
    c_{\mathrm{over}}(p-b), & p\geq b,\\
    c_{\mathrm{under}}(b-p), & p<b.
  \end{cases}
  \label{eq:asymmetric_identity}
\end{equation}
Hence, for a decision-loss budget $\eta\geq0$, the exact necessary-and-sufficient harm region is
\begin{equation}
  \mathcal H^{\mathrm{asym}}_{\eta}(b)
  =\left[b-\frac{\eta}{c_{\mathrm{under}}},
          b+\frac{\eta}{c_{\mathrm{over}}}\right].
  \label{eq:asymmetric_region}
\end{equation}
Replacing $I_{it}$ by
\begin{equation}
  I^{\mathrm{asym}}_{it}
  =\left[
  \max\!\left\{0,b_{it}-\frac{\eta}{c_{\mathrm{under}}}\right\},
  \min\!\left\{1.3,b_{it}+\frac{\eta}{c_{\mathrm{over}}},p_{i,t-1}\right\}
  \right]
  \label{eq:asymmetric_interval}
\end{equation}
preserves recursive nonemptiness, prefix causality, and non-increase when the costs and budget are constant within a trajectory. Setting $c_{\mathrm{under}}=c_{\mathrm{over}}=1$ and $\eta=\delta$ recovers PCHP exactly.
\end{corollary}
\begin{proof}
If $p\geq b$, the loss difference is bounded above by $c_{\mathrm{over}}(p-b)$ and attains that value for any $y\leq b$; if $p<b$, it is bounded above by $c_{\mathrm{under}}(b-p)$ and attains that value for any $y\geq b$. Solving the resulting two directional inequalities gives Eq.~\eqref{eq:asymmetric_region}. Because $b_{it}$ is non-increasing over $t$, the lower endpoint is likewise non-increasing over $t$, so the induction used in Theorem~\ref{thm:pchp} applies unchanged to Eq.~\eqref{eq:asymmetric_interval}.
\end{proof}

Directionally weighted absolute loss is an established forecasting primitive related to quantile loss \citep{koenker1978regression}; the contribution of Corollary~\ref{cor:asymmetric} is not the loss itself, but its exact outcome-agnostic harm geometry and recursively feasible integration into the same online contract.

\begin{corollary}[Exact viability under time-varying asymmetric contracts]
\label{cor:timevarying}
Let the budget $\eta_{it}\geq0$ and costs $c_{\mathrm{under},it},c_{\mathrm{over},it}>0$ vary with $t$, and define
\begin{align}
  L_{it}&=\max\!\left\{0,b_{it}-\frac{\eta_{it}}{c_{\mathrm{under},it}}\right\},
  &
  U_{it}&=\min\!\left\{1.3,b_{it}+\frac{\eta_{it}}{c_{\mathrm{over},it}}\right\},
  \label{eq:timevarying_bounds}\\
  I^{\mathrm{tv}}_{it}&=[L_{it},\min\{U_{it},p_{i,t-1}\}].
  \label{eq:timevarying_interval}
\end{align}
For a non-increasing protected state $b_{it}$:
\begin{enumerate}[label=(\roman*),leftmargin=*]
  \item given the realized prefix, $I^{\mathrm{tv}}_{it}$ is nonempty if and only if $L_{it}\leq p_{i,t-1}$;
  \item the recursion is feasible for every candidate sequence if and only if $L_{it}\leq L_{i,t-1}$ for every $t\geq2$;
  \item given $p_{i,t-1}$, the exact minimum current budget is
  \begin{equation}
    \eta^{\min}_{it}
    =c_{\mathrm{under},it}(b_{it}-p_{i,t-1})_{+}.
    \label{eq:minimum_viable_budget}
  \end{equation}
\end{enumerate}
The upper-cost schedule affects the upper tube radius and projected output, but not whether the recursive interval is empty.
\end{corollary}
\begin{proof}
Because $b_{it}\in[0,1.3]$, the standalone tube contains $b_{it}$ and hence $L_{it}\leq U_{it}$. Therefore Eq.~\eqref{eq:timevarying_interval} is nonempty exactly when its lower endpoint does not exceed $p_{i,t-1}$, proving (i). If $L_{it}\leq L_{i,t-1}$, every feasible preceding projection satisfies $p_{i,t-1}\geq L_{i,t-1}\geq L_{it}$, so induction proves sufficiency in (ii). Conversely, let $t^{\star}$ be the first index with $L_{i,t^{\star}}>L_{i,t^{\star}-1}$ and choose candidates at or below the lower endpoint through $t^{\star}-1$. The projection then emits $p_{i,t^{\star}-1}=L_{i,t^{\star}-1}$, making the next interval empty. Thus no weaker schedule condition can guarantee feasibility for every candidate sequence. Finally, solving $L_{it}\leq p_{i,t-1}$ for the smallest nonnegative $\eta_{it}$ gives Eq.~\eqref{eq:minimum_viable_budget}. The lower endpoint contains no $c_{\mathrm{over},it}$, which proves the final statement.
\end{proof}

For a predeclared schedule, Corollary~\ref{cor:timevarying} gives a full-trajectory certificate before execution. If costs or budgets arrive only prefix causally, Eq.~\eqref{eq:minimum_viable_budget} instead supplies an online viability diagnostic: a declared budget below the floor makes the interval empty and must trigger failure rather than silent budget enlargement.

\begin{corollary}[Exact constrained optimality]
\label{cor:optimality}
For every record, the PCHP output is the unique solution of
\begin{equation}
\begin{aligned}
  p_{it}=\operatorname*{arg\,min}_{p\in\R}\;&\frac{1}{2}(p-c_{it})^2\\
  \text{subject to}\quad
  &p\in\mathcal Y,\qquad p\leq p_{i,t-1},\\
  &\sup_{y\in\R}\bigl(|p-y|-|b_{it}-y|\bigr)\leq\delta.
\end{aligned}
\label{eq:constrained_optimality}
\end{equation}
\end{corollary}
\begin{proof}
Theorem~\ref{thm:harm_region} makes the feasible set in Eq.~\eqref{eq:constrained_optimality} exactly $I_{it}$. Theorem~\ref{thm:pchp} makes that interval nonempty. The strictly convex objective therefore has the unique minimizer $\Pi_{I_{it}}(c_{it})$, which is Eq.~\eqref{eq:method}.
\end{proof}

\begin{corollary}[Aggregation inheritance]
For any nonnegative weights $w_k$ and budgets $\delta_k$,
\begin{equation}
  \sum_k w_k\left(|p_k-y_k|-|b_k-y_k|\right)
  \leq \sum_k w_k\delta_k.
  \label{eq:aggregate}
\end{equation}
If $\sum_k w_k=1$, the weighted MAE increase is bounded by the weighted average budget.
\end{corollary}

No independence assumption is needed for Eq.~\eqref{eq:aggregate}. Statistical uncertainty for empirical utility remains separate and must respect the domain--cell--cycle hierarchy. The constant experimental budget is a recursively feasible special case; Corollary~\ref{cor:timevarying} states the exact additional condition required by a time-varying contract.

\subsection{Computational structure}

Once the two frozen regressors have produced the current record's predictions, PCHP stores only $(b_{i,t-1},p_{i,t-1})$ and performs a constant number of scalar operations. It therefore requires $O(1)$ state and $O(1)$ time per record after model inference. The bounded scalar squared-loss extension retains the same complexity because both endpoints in Eq.~\eqref{eq:squared_region} are closed form. The audit trail records $r_{it}$, $c_{it}$, $b_{it}$, $I_{it}$, $p_{it}$, $\alpha$, and $\delta$. Unlike retrospective isotonic regression, no trajectory sorting or future recomputation is required online.

\section{Experimental design}

\subsection{Development domains and nested model selection}

Development data were standardized from the BatteryLife processed snapshot (v11), BatteryML, Battery Archive, and the original experimental sources for CALCE, HNEI, HUST, MATR, MICH and MICH\_EXP, RWTH, SNL, UL\_PUR, SDU, and XJTU \citep{zhang2024batteryml,tan2025batterylife,batterylife2026processed,batteryarchive,xing2013ensemble,devie2018variability,ma2022personalized,severson2019cyclelife,attia2020closedloop,weng2021formation,li2021oneshot,preger2020degradation,juarezrobles2020safety,wang2025sorting,wang2024pinn}. After the first-five reference window, the corpus contained $12$ complete dataset domains, $586$ physical cells, and $601{,}932$ records (Table~\ref{tab:data}). All twelve domains had been opened during method development and are classified as internal or reused evidence rather than external evaluation.

\begin{table}[t]
\centering
\caption{Development corpus after the first-five reference window. Records are nested within physical cells, and cells are nested within complete dataset domains.}
\label{tab:data}
\begin{tabular}{lrr@{\qquad}lrr}
\toprule
Domain & Cells & Records & Domain & Cells & Records \\
\midrule
CALB & 27 & 2,629 & MICH\_EXP & 18 & 6,455 \\
CALCE & 13 & 14,233 & RWTH & 48 & 105,140 \\
HNEI & 14 & 15,085 & SDU & 86 & 48,716 \\
HUST & 77 & 145,735 & SNL & 61 & 96,092 \\
MATR & 169 & 139,148 & UL\_PUR & 10 & 2,195 \\
MICH & 40 & 19,681 & XJTU & 23 & 6,823 \\
\midrule
\multicolumn{4}{r}{Total} & 586 & 601,932 \\
\bottomrule
\end{tabular}
\end{table}

The outer evaluation held out one complete target domain. The two regressors were fitted on the remaining eleven domains. Within those source domains, each candidate value
\[
\alpha\in\{1,0.5,0.2,0.1,0.05,0.03,0.02,0.015,0.01,0.005\}
\]
was evaluated by holding out each source domain in turn and fitting on the other ten. The selected $\alpha$ minimized inner domain-equal cell-macro MAE; ties within $10^{-12}$ favored the larger $\alpha$. The outer target's labels never entered this selection. The harm budget was fixed at $\delta=0.01$ SOH units for all causal analyses. Because this budget was chosen during development, its decision meaning must be supplied by a deployment owner rather than inferred from the external stress result.

For each cell, MAE was first averaged across its post-reference records; cells were then averaged equally within a domain, and the $12$ domain estimates were averaged equally. The primary development estimand was the domain-equal paired difference between PCHP and its $\alpha$-matched protected causal baseline. This hierarchy prevents large archives or long-lived cells from dominating the cross-domain claim.

\subsection{Loss-geometry falsification protocol}

Before integrating the loss extension into the manuscript, we froze its symbolic identities, finite-interval formulas, exterior-point failure checks, recursive-feasibility gate, tolerance $10^{-12}$, and random seed $20{,}260{,}802$. The deterministic test read no battery outcome. It compared the bounded squared-loss closed form against endpoint maximization in $200{,}000$ random configurations, tested both feasible boundaries and adjacent exterior points, and projected arbitrary candidates along $2{,}000$ random non-increasing protected trajectories of length $128$. Separate counterexamples moved the hidden outcome toward the appropriate infinity for nontrivial upward and downward squared-loss updates. Any formula mismatch, feasible-point violation, nonempty exterior point, empty recursive interval, output increase, or failure to diverge required rejection without tolerance adjustment.

\subsection{Operator-stability falsification protocol}

Before integrating Theorem~\ref{thm:nonexpansive}, we froze a label-free validator with seed $20{,}260{,}802$, absolute tolerance $2\times10^{-12}$, exact local grids, and $30{,}000$ independently generated trajectory pairs. The validator checked $3{,}888$ local protected-state configurations, $21{,}168$ local projection configurations, $108$ proximal equivalences, both prefix inequalities at every simulated step, and an explicit tightness witness. Any positive excess beyond tolerance, proximal mismatch, or failure to attain constant one required rejection without changing the theorem or tolerance.

\subsection{Mechanism falsification and distinguishing controls}

Before nested scoring, the smallest mechanism test used label-free development predictions only. At prefix fractions $0.25$, $0.50$, and $0.75$, we compared outputs computed on the prefix with the same positions recomputed after the full future trajectory became available. A change above $10^{-4}$ SOH units counted as a retrospective revision. The causal mechanism was required to have zero prefix difference, while full-trajectory isotonic processing was required to revise at least half the cells in at least nine domains. A second intervention added a $+0.05$ or $-0.05$ innovation to one protected raw prediction. The predicted signature was exact rejection of the positive innovation, immediate absorption of the negative innovation at rate $\alpha$, and reduced cumulative disturbance for $\alpha=0.02$ relative to $\alpha=1$ in at least $95\%$ of interventions in at least nine domains.

The principal comparator was the causal protected baseline generated with the same source-selected $\alpha$. A retrospective isotonic harm-capped estimator was retained only as a leakage audit: it has access to the complete target trajectory and is therefore not deployable under the stated information boundary. The strongest alternative explanation was that PCHP merely shifts the protected state by a constant inside the same $\delta$-tube. To test it, a candidate-free control jointly selected $\alpha$ and a constant offset from $\{-0.01,-0.009,\ldots,0.01\}$ using only the eleven source domains of each outer fold. It obeyed the same range, monotonicity, and harm-budget constraints as PCHP.

After the development outcomes had been opened, we protocol-locked a retrospective budget-path sensitivity to test whether this distinction depended on the single value $\delta=0.01$. Before computing the path, we fixed
\[
\delta\in\{0,0.0025,0.005,0.0075,0.01,0.015,0.02,0.03\},
\]
the original $\alpha$ grid, the outer-domain split, one integrated primary estimand, bootstrap seed $20{,}260{,}811$, and decision gates. For every outer domain, budget, and $\alpha$, the candidate-free comparator used the weighted-median source residual---the exact constant offset minimizing inner domain-equal cell-macro absolute error---clipped to $[-\delta,\delta]$; $\alpha$ was then selected on the same eleven source domains. No outer-target label entered either selection. The zero-budget arm was required to reproduce the protected state, and all arms were required to satisfy the declared displacement, regret, range, and monotonicity checks. This analysis is protocol-locked but retrospective, not a new independent evaluation.

We also constructed a stronger \emph{unprotected} causal comparator to measure the opportunity cost of the tube. For each outer target, the same source-only double-holdout design selected $\alpha$ for
\[
q_{i1}=\Pi_{[0,1.3]}(c_{i1}),\qquad
q_{it}=\Pi_{[0,1.3]}\!\left(q_{i,t-1}+\alpha\min(c_{it}-q_{i,t-1},0)\right).
\]
This comparator is prefix causal, range bounded, and non-increasing, but it has no protected-baseline tube. All twelve outer folds selected $\alpha=0.05$ from source domains only. It is therefore an operationally strong accuracy boundary, not a harm-budget-matched baseline.

\subsection{Outcome-blind, prospectively locked BaSyTec external validation}

The Bayreuth degradation-path dataset contains $48$ commercial lithium-ion cells cycled for $70$ ageing cycles under $24$ temperature--C-rate combinations, with two nominal replicates per condition \citep{ramasubramanian2025degradation,roder2025degradationdata}. The public Zenodo record, all $48$ cell archives, and their official checksums were snapshotted before scoring. The first parser specification failed closed before any capacity value was read because the raw export contained cumulative \texttt{Ah[Ah]} rather than a per-cycle capacity field. We preserved that failure. Before inspecting any cumulative-Ah value, cell F0001 was permanently designated as the parser-development cell and excluded from evaluation. A subsequent header-only audit read zero data rows across the remaining $47$ archives and froze the two observed delimiter families and four header signatures before prediction.

For consecutive rows $j-1,j$ within cycle $c$, cumulative ampere-hours $A_j$ and current $I_j$ defined the frozen discharge-capacity reconstruction
\begin{equation}
Q_c=\sum_{j=2}^{n_c}\mathbf{1}\!\left(I_j<-10^{-5}\,\mathrm{A}\right)
\max(A_{j-1}-A_j,0).
\label{eq:basytec_capacity}
\end{equation}
The label-blind stage selected only cycle identity plus charge-side time, current, and voltage; it converted numeric hours to seconds, extracted the fixed $600\,\mathrm{s}$ charge window, and hashed $3{,}101$ post-anchor predictions for $47$ cells before Eq.~\eqref{eq:basytec_capacity} was evaluated. The scorer then formed Eq.~\eqref{eq:soh}. A cell required a first-five median capacity in $[0.08,0.16]\,\mathrm{Ah}$ and at least ten post-anchor labels. The frozen PCHP settings were $\alpha=0.01$ and $\delta=0.01$; the unprotected boundary used the source-selected $\alpha=0.05$.

The decision gate required at least $36$ eligible cells, a negative cell-equal paired MAE difference, improvement in at least two thirds of non-tied cells, a cell-bootstrap $95\%$ interval with upper endpoint below zero, and maximum pointwise displacement, observed absolute-loss regret, and cell-macro harm no greater than $0.01$ within $10^{-12}$ tolerance. The remaining $47$ cells were held out from method, hyperparameter, estimand, exclusion-rule, and gate development; the available workspace records their prediction hashes before capacity access. We therefore treat the analysis as outcome-blind, prospectively locked external validation on held-out public data. The label describes the frozen evaluation sequence: the archive was previously collected and public, so this is not a prospectively collected laboratory study or an independent-laboratory replication.

\subsection{Protocol-frozen NASA one-shot stress test}

NASA PCoE Battery Data Set~5 was selected from official metadata and frozen before archive download \citep{saha2007nasa}. The protocol fixed the official archive identity, inclusion of all recursively discovered MATLAB files, charge--discharge pairing, the $600\,\mathrm{s}$ charge window, first-five reference, feature fields, $\alpha=0.01$, $\delta=0.01$, minimum cell size, estimand, decision gates, and hashes of both prediction and scoring programs. The archive contained $38$ MATLAB files representing $34$ physical batteries after duplicate identities were reconciled.

Stage one accessed only cycle type and time plus voltage, current, temperature, and time arrays from charge records. It created and hashed $2{,}598$ post-anchor predictions for $34$ batteries without reading discharge capacity. Stage two independently read the frozen discharge locators and the scalar \texttt{Capacity} field, formed Eq.~\eqref{eq:soh}, and scored the frozen predictions once. A battery required at least ten valid post-anchor records. The frozen outcome contract yielded $33$ batteries and $2{,}556$ scored records; $44$ discharge records with a non-scalar, non-finite, or non-positive capacity were rejected by the prespecified rule. No post-outcome roster, reference, model, parameter, or endpoint change was permitted.

The external gate required a negative cell-equal paired MAE difference, improvement in at least $75\%$ of eligible batteries, maximum cell-macro harm no greater than $0.01$, and all deterministic certificates. The archive is one repository identity from one laboratory program, so batteries are independent units for within-dataset replication but not independent laboratories. In the available project lineage, the frozen prediction artifact preceded access to Dataset~5 discharge capacities. We therefore describe this as a protocol-frozen external stress test; because its absolute target is incompatible, it does not serve as positive external utility confirmation or as a prospectively collected laboratory study.

\subsection{Statistical analysis}

The complete dataset domain was the primary analysis unit for development effects; the physical cell was the unit for the BaSyTec and NASA within-program effects. Repeated records were nested and never counted as independent replications. Development and candidate-control intervals used $100{,}000$ percentile bootstrap replicates of the $12$ paired domain effects with fixed seeds $20{,}260{,}802$ and $20{,}260{,}801$, respectively. The three secondary-backbone comparisons formed one prespecified robustness family and used Bonferroni-adjusted two-sided $98.33\%$ percentile intervals from $100{,}000$ complete-domain resamples with seed $20{,}260{,}802$. For the budget path, the prespecified domain-level estimand was
\begin{equation}
A_d=\frac{1}{0.03}\int_{0}^{0.03}
\left[\operatorname{MAE}^{\mathrm{PCHP}}_d(\delta)-
\operatorname{MAE}^{\mathrm{shift}}_d(\delta)\right],\mathrm d\delta,
\label{eq:budget_auc}
\end{equation}
evaluated by the trapezoidal rule on the locked grid. The primary report is the domain-equal mean of the $12$ values $A_d$ and its $100{,}000$-replicate domain-bootstrap interval; the per-budget intervals are descriptive slices of this one path, not separate confirmation claims. Because MICH and MICH\_EXP originate from the same experimental program, a conservative post-hoc sensitivity merged them before source-group-equal aggregation and bootstrap resampling; it also recomputed the mean after deleting each source group in turn. BaSyTec uncertainty used $100{,}000$ paired-cell bootstrap replicates with seed $20{,}260{,}808$; a post-hoc sensitivity averaged the two cell replicates within each temperature--C-rate condition and resampled complete conditions. NASA uncertainty used $100{,}000$ paired-battery bootstrap replicates with seed $20{,}260{,}801$. The two-sided exact sign test excluded ties and was descriptive rather than a confirmation gate. All development domains had been opened, so their intervals summarize heterogeneity over the observed domain set rather than provide confirmatory coverage for a prespecified superpopulation. The two public-data workflows used frozen multi-criterion decision gates rather than alpha-calibrated hypothesis tests. Outside the explicitly adjusted three-backbone robustness family, no familywise type-I-error claim or multiplicity-adjusted inference is made. Each interval corresponds to its stated comparison, and no cycle-level test or pooled-row confidence interval was used. Deterministic certificates are pointwise mathematical properties and do not depend on these resampling assumptions.

Because only $12$ complete domains support the development estimands, a clearly post-hoc finite-sample sensitivity additionally enumerated all $2^{12}=4{,}096$ paired sign assignments for each of the primary, candidate-control, and integrated budget-path mean effects. It also computed Student-$t$ intervals, exact paired Wilcoxon tests, and exact direction-only sign tests. These complementary calculations were not used for method selection, do not replace the frozen percentile-bootstrap reports, and do not convert opened development evidence into confirmation.

\section{Results}

\subsection{The loss-dependent geometry passes symbolic and recursive falsification}

All $27$ frozen aggregate checks passed. Across the $200{,}000$ bounded squared-loss configurations, the maximum absolute discrepancy between endpoint maximization and Eq.~\eqref{eq:squared_identity} was $1.33\times10^{-15}$; the active lower and upper boundary discrepancies from the declared budget were at most $8.88\times10^{-16}$ and $1.55\times10^{-15}$, respectively. Across $256{,}000$ recursive records, no interval was empty, the maximum output increase was zero, and the maximum budget excess was $7.77\times10^{-16}$. Both unbounded-outcome constructions diverged in the direction predicted by Eq.~\eqref{eq:squared_unbounded}. These tests do not replace the proofs; they exclude sign, endpoint, clipping, and recursion errors in the implemented formulas.

\subsection{The causal operator prevents future revision and has the predicted shock response}

The smallest falsification passed both gates. Full-trajectory isotonic processing revised more than half of the evaluated cell prefixes in all $12$ domains, whereas the maximum difference between a causal prefix and the same prefix inside any future extension was exactly $0$. Under the controlled $+0.05$ intervention, the maximum disturbance of the one-sided protected state was $0$. Under the $-0.05$ intervention, the maximum error from the predicted immediate response $\alpha\times0.05$ was at most $2.00\times10^{-16}$ in the saved intervention table, i.e., within floating-point tolerance. Relative to $\alpha=1$, $\alpha=0.02$ reduced cumulative shock disturbance for at least $95\%$ of interventions in every domain. Figure~\ref{fig:causality} shows that the mechanism is an information-boundary change, not a relabelled offline smoother.

\begin{figure}[t]
  \centering
  \includegraphics[width=\textwidth]{fig3_prefix_causality.pdf}
  \caption{Mechanism falsification across $12$ development domains and $586$ physical cells. \textbf{A}, fraction of cells whose historical offline output changed by more than $10^{-4}$ after future records were appended; PCHP's maximum prefix difference was $0$. The dashed line marks the prespecified one-half-cell gate. \textbf{B}, paired domain medians of cumulative shock disturbance under $\alpha=1$ and $\alpha=0.02$ on a logarithmic scale; every intervention improved in every domain. Records and interventions are nested within cells; domains are the cross-domain evaluation units.}
  \label{fig:causality}
\end{figure}

\subsection{Nested source-only selection improves cross-domain utility within the harm budget}

Under strict nested leave-one-domain-out evaluation, PCHP changed domain-equal cell-macro MAE from $0.07524$ for its protected causal baseline to $0.07038$, a $6.46\%$ relative reduction. The paired difference was $-0.00486$ with a domain-cluster $95\%$ bootstrap interval of $[-0.00671,-0.00301]$; $11$ domains improved, none tied, and one worsened. The largest domain-level harm was $0.00109$, well inside the $0.01$ pointwise budget. Figure~\ref{fig:domain_effects} reports every domain rather than presenting the $601{,}932$ repeated records as independent evidence.

\begin{figure}[t]
  \centering
  \includegraphics[width=0.92\textwidth]{fig2_domain_effects.pdf}
  \caption{Paired PCHP-minus-baseline cell-macro MAE for all $12$ held-out development domains. Each point is one complete dataset domain after cell-equal aggregation; negative values favor PCHP. The diamond and horizontal interval denote the domain-equal mean and its $95\%$ percentile bootstrap interval over $100{,}000$ resamples of complete domains.}
  \label{fig:domain_effects}
\end{figure}

The causal method also had a lower point estimate than the retrospective harm-capped estimator: $0.07038$ versus $0.07232$. Their paired difference was $-0.00195$ with a $95\%$ interval of $[-0.01073,0.00606]$. This comparison does not establish statistical superiority; it indicates that enforcing the deployable information boundary did not produce a clear domain-level accuracy penalty in the opened development corpus.

\subsection{The output contract transfers across untuned secondary backbones}

To test whether the observed utility was specific to ExtraTrees, the protocol-locked retrospective audit replaced both frozen regressors with Ridge, HistGradientBoosting, or LightGBM and reused the exact outer-domain $\alpha$ schedule selected for the primary backbone. No secondary-backbone tuning or target-domain label access was allowed. For Ridge, HistGradientBoosting, and LightGBM, respectively, PCHP-minus-protected-state domain-equal cell-macro MAE differences were $-0.00419$, $-0.00425$, and $-0.00422$. Their Bonferroni-adjusted $98.33\%$ intervals, in the same order, were
\[
[-0.00677,-0.00131],\qquad
[-0.00633,-0.00203],\qquad
[-0.00621,-0.00231].
\]
Domain wins/ties/losses were $10/0/2$, $11/0/1$, and $11/0/1$. All displacement, observed-regret, range, monotonicity, and exact prefix-replay checks passed. This result supports portability across these three untuned families and excludes a narrow ExtraTrees-specific explanation; it does not establish universal model independence, deep-model portability, or best-backbone superiority.

\subsection{The learned candidate contributes information beyond a constant budget-feasible shift}

The candidate-free control attained domain-equal cell-macro MAE $0.07515$, compared with $0.07038$ for PCHP, a $6.35\%$ relative reduction. Expressed as PCHP minus control, the paired difference was $-0.00477$ with a $95\%$ domain-bootstrap interval of $[-0.00808,-0.00159]$. PCHP improved $9$ domains and worsened $3$, with maximum harm relative to this control of $0.00424$.

The protocol-locked retrospective path made the control exact rather than grid approximated and repeated the comparison from zero to three SOH percentage points. All seven positive budgets favored PCHP in the domain-equal mean; the differences ranged from $-0.00189$ at $\delta=0.0025$ to $-0.01205$ at $\delta=0.03$, and every descriptive per-budget interval lay below zero. The primary budget-normalized AUC difference in Eq.~\eqref{eq:budget_auc} was $-0.00694$, with a domain-bootstrap interval of $[-0.01108,-0.00296]$ and domain wins/ties/losses of $11/0/1$. The zero-budget identity, displacement, regret, range, and monotonicity gates all passed. An independent hash-pinned audit retained the original failed validator record but showed that its sole failed anchor compared the high-precision value $0.0703775121469438$ with a ten-decimal JSON serialization under a $10^{-12}$ tolerance; the authoritative CSV aggregates, all $12$ domain values, and the source-only selection roster matched the historical run. Figure~\ref{fig:candidate_control} therefore weakens the single-budget fixed-bias explanation while preserving the analysis boundary: development outcomes had already been opened, so this path is retrospective robustness evidence rather than independent evidence.

The source-group sensitivity did not rely on treating MICH and MICH\_EXP as independent programs. After merging them, the primary source-group-equal difference was $-0.00469$ with a $95\%$ bootstrap interval of $[-0.00667,-0.00274]$ and wins/ties/losses of $10/0/1$. The candidate-control difference was $-0.00467$ with interval $[-0.00825,-0.00125]$ and wins/ties/losses of $8/0/3$. For the primary and control comparisons, respectively, all leave-one-source-group-out means remained negative within $[-0.00526,-0.00417]$ and $[-0.00556,-0.00347]$. This post-hoc sensitivity addresses shared provenance; it does not replace the frozen complete-domain estimand.

\begin{figure}[t]
  \centering
  \includegraphics[width=\textwidth]{fig4_candidate_control.pdf}
  \caption{Protocol-locked retrospective budget-path control for candidate information. \textbf{A}, domain-equal PCHP-minus-exact-budget-feasible-shift cell-macro MAE over the locked budget grid; negative values favor PCHP, vertical intervals are descriptive domain-bootstrap intervals, and the diamond marks the manuscript setting $\delta=0.01$. For every budget and outer fold, the control uses no candidate features and selects $\alpha$ plus the exact source-optimal constant budget-feasible shift using only the other eleven domains. \textbf{B}, the $12$ domain-specific budget-normalized AUC differences from Eq.~\eqref{eq:budget_auc}; the diamond and interval are the domain-equal mean and its $95\%$ percentile bootstrap interval from $100{,}000$ complete-domain resamples. The development outcomes had been opened before this protocol was locked, so the figure reports retrospective mechanism sensitivity, not independent evidence.}
  \label{fig:candidate_control}
\end{figure}

\subsection{A strong unprotected comparator exposes the harm-control--accuracy opportunity cost}

The source-tuned unprotected causal comparator reduced domain-equal cell-macro MAE from $0.07524$ to $0.05470$. Its paired difference from the protected baseline was $-0.02054$ with a $95\%$ domain-bootstrap interval of $[-0.03364,-0.00675]$ and wins/ties/losses of $10/0/2$. PCHP retained $23.7\%$ of this gain: its MAE exceeded the unprotected comparator by $0.01568$, with interval $[0.00350,0.02725]$ and wins/ties/losses of $3/0/9$. After merging MICH and MICH\_EXP, utility retention was $24.9\%$ and the same qualitative conclusion held.

This additional accuracy was not safety matched. The unprotected comparator exceeded the $\delta=0.01$ displacement budget in all $12$ domains, all $586$ cells, and all $524{,}165$ scored records, with a maximum displacement of $0.35680$. PCHP had no budget violation. The comparison therefore falsifies a dominance claim: PCHP does not beat the strongest causal candidate on accuracy. Its contribution is to occupy a declared worst-case harm region while retaining a measurable fraction of the available utility.

\subsection{BaSyTec evaluation passes under the frozen contract}

The BaSyTec gate returned \texttt{RETAIN}. Of the $47$ evaluation cells, $45$ satisfied the frozen reference-capacity range, yielding $2{,}969$ scored cycles. Cell-equal MAE changed from $0.17496$ for the protected causal baseline to $0.16510$ for PCHP. The paired difference was $-0.00986$; its paired-cell bootstrap $95\%$ interval extended from $-0.01000$ to $-0.00959$, and all $45$ eligible cells improved. Independent reconstruction matched every cell aggregate within $10^{-12}$. A post-hoc condition-equal sensitivity was negative in all $23$ represented temperature--C-rate conditions, with an interval from $-0.01000$ to $-0.00958$. Maximum PCHP displacement and maximum observed absolute-loss regret were $0.01$ within floating-point tolerance, while the maximum cell-macro difference was $-0.00465$.

Both cells at $0\,^{\circ}\mathrm{C}$ and $1.5C$ were excluded by the prespecified $0.08\,\mathrm{Ah}$ lower reference bound: F0009 and F0048 had first-five medians of $0.06862$ and $0.06935\,\mathrm{Ah}$. A clearly post-hoc sensitivity that included them preserved the result across all $47$ cells and all $24$ conditions: the paired difference was $-0.00986$, its interval was $[-0.01000,-0.00961]$, and $47/47$ cells improved. This sensitivity shows that the exclusion did not create the effect, but it does not replace the frozen $45$-cell estimand.

The same evaluation exposes why the empirical claim must remain narrow. The source-tuned unprotected comparator attained MAE $0.06172$, but its maximum displacement was $0.28508$ and it violated the budget in every eligible cell and all $2{,}969$ records; PCHP retained only $8.7\%$ of its accuracy gain. Moreover, the protected baseline underpredicted every scored record, PCHP saturated the positive $0.01$ boundary in $99.1\%$ of records, and a post-hoc fixed $+0.01$ shift achieved MAE $0.16496$, marginally below PCHP by $0.00014$. Thus, Figure~\ref{fig:basytec} evaluates the frozen harm-and-utility contract and quantifies its opportunity cost; this external dataset alone does not distinguish adaptive PCHP behavior from a simple budget-feasible bias correction.

\begin{figure}[t]
  \centering
  \includegraphics[width=\textwidth]{fig6_basytec_confirmation.pdf}
  \caption{Outcome-blind, prospectively locked BaSyTec external validation and harm-control--accuracy boundary. \textbf{A}, paired PCHP-minus-baseline MAE for all $45$ eligible cells; negative values favor PCHP. The diamond and interval are the cell-equal mean and $95\%$ paired-cell bootstrap interval from $100{,}000$ resamples. \textbf{B}, cell-equal MAE against maximum pointwise displacement from the protected baseline. PCHP remains inside the declared $\delta=0.01$ region. The fixed $+0.01$ shift is a post-hoc diagnostic, whereas the source-tuned unprotected comparator is prefix causal but violates all eligible-cell budgets. Cycles are nested within cells.}
  \label{fig:basytec}
\end{figure}

\subsection{The frozen NASA stress test retains relative utility but exposes target incompatibility}

The one-shot NASA gate returned \texttt{RETAIN}. Across $33$ eligible physical batteries, cell-equal MAE changed from $0.73379$ for the causal protected baseline to $0.72791$ for PCHP. The paired difference was $-0.00588$ with a battery-bootstrap $95\%$ interval of $[-0.00784,-0.00373]$; $27$ batteries improved and $6$ worsened. The maximum battery-level harm was $0.00849$. The two-sided exact sign-test value was $0.000324$, reported descriptively because it was not a frozen gate. Every prediction satisfied $|p_{it}-b_{it}|\leq0.01$, the maximum observed absolute-loss increase was $0.01$ within floating-point tolerance, and all protected and final trajectories were non-increasing.

The absolute errors do not support a high-accuracy NASA SOH claim. After the one-shot release, the frozen label audit found that some discharge capacities were not comparable under the first-five reference. Battery B0041 had a reference capacity of $0.05584\,\mathrm{Ah}$ and a maximum normalized value of $21.77$; B0039 had a reference of $0.47114\,\mathrm{Ah}$ and a maximum normalized value of $3.76$. Across released labels, the $95$th percentile was $1.347$ and the maximum was $21.77$, whereas model outputs were frozen to $[0,1.3]$. Figure~\ref{fig:nasa} therefore presents the relative effect and the protocol boundary together. Removing these batteries or redefining the reference after outcome access would invalidate the frozen experiment and was not done.

\begin{figure}[t]
  \centering
  \includegraphics[width=\textwidth]{fig5_nasa_stress.pdf}
  \caption{Protocol-frozen NASA stress test. \textbf{A}, paired PCHP-minus-baseline MAE for all $33$ eligible physical batteries; negative values favor PCHP, and the diamond with interval denotes the cell-equal mean and $95\%$ paired-battery bootstrap interval from $100{,}000$ resamples. \textbf{B}, released normalized capacity labels by battery on a logarithmic axis, with the frozen model range $[0,1.3]$ indicated. The same display shows why the relative harm-control result is valid while a high-accuracy absolute SOH claim is not. The $2{,}556$ records are repeated measurements nested within batteries.}
  \label{fig:nasa}
\end{figure}

\begin{table}[t]
\centering
\caption{Key paired effects. Differences are compared-method minus reference MAE; intervals resample the stated independent unit, not individual records. The unprotected row is an accuracy boundary and violates the harm budget.}
\label{tab:primary}
\small
\begin{tabularx}{\textwidth}{>{\raggedright\arraybackslash}Xrrrrl}
\toprule
Evidence tier and comparison & Unit $n$ & Reference MAE & Compared MAE & Difference & $95\%$ interval \\
\midrule
Nested development & 12 domains & 0.07524 & 0.07038 & $-0.00486$ & $[-0.00671,-0.00301]$ \\
PCHP vs candidate-free budget-feasible shift & 12 domains & 0.07515 & 0.07038 & $-0.00477$ & $[-0.00808,-0.00159]$ \\
Unprotected candidate vs baseline & 12 domains & 0.07524 & 0.05470 & $-0.02054$ & $[-0.03364,-0.00675]$ \\
BaSyTec PCHP vs baseline & 45 cells & 0.17496 & 0.16510 & $-0.00986$ & $[-0.01000,-0.00959]$ \\
NASA one-shot stress & 33 batteries & 0.73379 & 0.72791 & $-0.00588$ & $[-0.00784,-0.00373]$ \\
\bottomrule
\end{tabularx}
\end{table}

\section{Discussion}

\subsection{What capability changes}

PCHP changes the operational contract of cross-domain battery prediction. A learned model remains responsible for proposing a useful update, but it no longer has unilateral authority over a deployed record. The protected causal state defines the prediction to be preserved, and Corollary~\ref{cor:viability_kernel} identifies the recursive feasible interval as the exact set of current outputs that remain extendable under every future non-increasing protected-state path. Corollary~\ref{cor:optimality} then shows that the final output is not an arbitrary point in that kernel: it is the unique update that preserves maximum fidelity to the current learned proposal under the complete declared constraint set. Theorem~\ref{thm:nonexpansive} adds that the complete recursive cascade cannot magnify a bounded discrepancy in either frozen predictor when the output contract is fixed. The resulting guarantees are simultaneously outcome agnostic and executable online. This is different from retrospective smoothing, which may be label free yet still violates the future-information boundary, and from unconstrained target adaptation, whose gain can become negative in an unknown domain.

The secondary-backbone audit adds an empirical boundary to that structural claim: the same frozen output contract remained useful for one linear and two boosting families without backbone-specific retuning. It rules out ExtraTrees-specific utility within the audited families, not dependence on every possible candidate architecture.

Theorem~\ref{thm:harm_region} also prevents a misleading interpretation of ``safe'' updating. Without assumptions on the hidden outcome, no method can move a point prediction and remain no worse than its baseline for every possible truth. A positive budget is therefore not a proof relaxation introduced for convenience; it is the unavoidable price of allowing any nontrivial outcome-free update. PCHP makes that price explicit and then spends it only inside the exact intersection with the online trajectory constraints.

The candidate-control result is important for interpreting the method as more than conservative clipping. If a source-selected constant shift had matched PCHP, the learned candidate would have been unnecessary and the contribution would reduce to a generic harm tube. Instead, commissioning-referenced, record-specific information contributes utility after the budget-feasible shift and assimilation rate have been matched. The integrated path further shows that this distinction is not an artifact of the single manuscript budget: the control is source-optimal among constant budget-feasible shifts at each locked budget, yet the domain-equal advantage remains throughout the positive grid. The data do not identify an electrochemical causal mechanism; they identify the operational value of candidate-specific information under a controlled output contract. Nor does the path calibrate an engineering budget: it was run retrospectively on opened development domains. The distinction is not universal either; BaSyTec exhibits nearly uniform baseline underprediction, under which a post-hoc fixed positive shift marginally outperforms PCHP. We therefore claim adaptive value from the nested development control and its protocol-locked sensitivity, not from every external dataset.

\subsection{Battery-specific operational meaning}

The projection theory can be reused outside batteries, but the present contract is not attached to the battery application after the fact. Battery SOH creates a specific information pattern: partial charge measurements can arrive during ordinary operation, while reference capacity outcomes require operating conditions that may be delayed or deliberately scheduled \citep{liu2023rapid,zhang2024flexible}; early records can be informative before substantial capacity degradation becomes visible \citep{severson2019cyclelife}; and a new supplier, chemistry, or protocol can shift the relation between charging response and capacity. These conditions determine the filtration, the commissioning reference, the cross-domain candidate, and the outcome-free harm requirement. Long-term capacity fade motivates the non-increasing emitted trend, but that constraint is an auditable engineering representation rather than an electrochemical law, which is why raw observations and rebound diagnostics remain necessary. The energy contribution is therefore an executable rule for admitting learned battery-health updates before reference capacity becomes available, with a declared cost in prediction loss. Applying the same mathematics to another asset would require an analogous delayed-outcome and protected-record contract.

\subsection{Meaning of the BaSyTec evaluation and the harm-control--accuracy boundary}

The BaSyTec experiment tests whether a prediction artifact frozen before capacity access improves its protected baseline while preserving all declared harm checks under a new laboratory, cell type, and $24$-condition stress grid. It passes that test, and the post-hoc full-condition sensitivity shows that the prespecified low-reference exclusion did not generate the direction of effect. This is stronger utility evidence than the NASA stress test because the resulting normalized capacities remain within a plausible SOH regime for the scored cells.

The same experiment prevents an inflated interpretation. Almost every PCHP output reaches the positive edge of the tube, and a fixed positive shift performs marginally better. Consequently, the result supports transfer of the harm-control and relative-utility contract, but not external identification of the adaptive candidate mechanism. The strong unprotected comparator sharpens a different boundary: much more accuracy was available, but only by leaving the declared harm tube for every cell. This is not a defect hidden by the guarantee; it is the engineering trade-off that the guarantee makes measurable. Whether $8.7\%$ retained utility is acceptable depends on decision costs that the present public data do not identify.

\subsection{Meaning of the NASA result}

The NASA experiment tests whether a completely frozen causal update preserves relative utility and deterministic harm control after a substantial data-source shift. It passes that test. It does not test whether the frozen source model produces deployment-grade absolute SOH under every NASA discharge regime. The unusually large normalized capacities reveal that the chosen target definition loses comparability when the first-five discharge records and later records use materially different conditions. The correct inference is therefore bounded: the update improved the protected baseline within its declared budget under a severe protocol shift, while the protected baseline and label contract remained inadequate for absolute deployment.

This distinction illustrates why deterministic and empirical claims must remain separate. Equation~\eqref{eq:pointwise_harm} remains true even when the baseline is poor, because it compares two predictions at the same record. Accuracy, calibration, and operational usefulness require compatible targets and representative deployment data; no output-space theorem can supply those conditions.

\subsection{Choosing the harm budget}

The budget $\delta$ is expressed in SOH units and should be tied to the consequence of changing a protected record. A smaller value preserves the baseline more tightly but limits candidate utility; a larger value admits more gain and more worst-case absolute-loss increase. The present $\delta=0.01$ is a frozen research setting, not a universal engineering tolerance. When underestimation and overestimation have unequal consequences, Corollary~\ref{cor:asymmetric} uses their positive cost weights and a budget $\eta$ in decision-loss units to create different allowable downward and upward radii. Only the cost ratio and the jointly scaled budget determine the standalone interval, so an operator must elicit and document this contract rather than infer it from the present accuracy experiments. Under a time-varying contract, Eq.~\eqref{eq:minimum_viable_budget} additionally reports the least budget that preserves the already emitted prefix; it diagnoses an infeasible declaration but does not calibrate the operational cost. Threshold or nonlinear decision losses still require a new feasible-set derivation and proof.

\subsection{Statistical interpretation}

The deterministic certificate is pointwise and requires no independent sampling. Empirical utility uncertainty has a different scope. In development, $12$ complete dataset domains support a cross-domain effect; $601{,}932$ records do not create that many independent transfers. The source-group sensitivity further shows that the result is not driven by counting two subsets of one experimental program separately. In BaSyTec, $45$ eligible cells support the frozen within-program effect and $23$ represented conditions support a post-hoc condition-cluster sensitivity; the full $47$-cell, $24$-condition analysis is explicitly post hoc. In NASA, $33$ batteries support within-program replication. Neither public workflow creates independent laboratories. The reported intervals therefore resample domains, conservative source groups, cells, conditions, or batteries as dictated by the stated analysis. They quantify heterogeneity over the observed unit set and do not establish universality across all chemistries, laboratories, temperatures, or duty cycles.

The finite-sample sensitivity preserved the negative mean direction without relying on the percentile bootstrap alone. For the primary effect and the integrated budget path, exact sign-flip values were $0.00098$ and $0.00781$, and Student-$t$ $95\%$ intervals were $[-0.00703,-0.00269]$ and $[-0.01172,-0.00217]$, respectively. For PCHP versus the candidate-free shift, the corresponding values were $0.01660$ and $[-0.00858,-0.00097]$. However, its direction-only sign-test value was $0.145996$ because PCHP won in $9/12$ domains. The supported claim is therefore a domain-equal mean utility contribution with heterogeneous domain effects, not universal dominance or a claim that a randomly drawn new domain is more likely than not to improve. All values are descriptive post-hoc sensitivities on opened development domains.

The confirmatory status of the BaSyTec analysis follows from design independence, not from requiring newly collected data or an ever larger number of archives. Its capacity outcomes did not determine the method, fitted predictors, hyperparameters, estimand, exclusion rule, or decision gate, and the scored predictions were fixed before those outcomes were accessed. Public availability and single-program provenance narrow the claim to an outcome-blind, protocol-locked within-program external confirmation; they do not turn the analysis into development evidence. Conversely, adding more previously opened datasets would increase breadth but would not create a stronger confirmatory boundary. Independent laboratory replication remains valuable for wider transport claims, but it is not required for the scoped empirical claim made here.

\section{Limitations and failure modes}

First, the guarantee is relative to the specified protected causal state. PCHP cannot diagnose an inaccurate baseline, repair incompatible labels, establish calibration, or guarantee electrochemical safety. Theorem~\ref{thm:nonexpansive} is likewise an incremental algorithmic-stability statement: it does not bound baseline error, statistical uncertainty, distribution shift, or physical risk. Second, Theorem~\ref{thm:metric_harm} covers metric losses, while Theorem~\ref{thm:squared_geometry} establishes the bounded-interval squared-loss construction; it does not justify an incorrect or unverifiable outcome bound. The empirical study still evaluates the absolute-loss contract; it does not establish squared-loss utility. Corollary~\ref{cor:timevarying} covers positive time-varying directional weights and nonnegative budgets only when its lower-endpoint condition holds. Threshold, discontinuous, or other non-metric decision losses require a new exact feasible-set derivation, and the present work does not estimate deployment cost weights. Third, strict non-increase estimates a long-term trend and can suppress genuine short-term capacity regeneration or measurement rebound; both raw observations and trend outputs should therefore be retained.

Fourth, all $12$ development domains were previously opened. The nested design prevents current outer-target labels from selecting $\alpha$, and the budget-path protocol was locked before its computation, but both analyses remain retrospective development evidence. Fifth, the BaSyTec workflow required two documented structural restarts before outcome access: one parser-development cell was excluded, and a header-only audit of the $47$ evaluation files preceded the final freeze. Neither step used the held-out cells' capacity outcomes, and neither changed the method, hyperparameters, estimand, exclusion logic, or decision gates. Predictions for those cells were fixed before capacity access; the $45$-cell analysis therefore qualifies as outcome-blind, prospectively locked external validation. Because the archive was already public and came from one laboratory program, the study is not prospectively collected and does not constitute independent-laboratory replication. Sixth, the frozen BaSyTec reference range excluded both cells from the $0\,^{\circ}\mathrm{C}$, $1.5C$ condition; the full-condition sensitivity is reassuring but post hoc. Seventh, the NASA archive is one public laboratory program with heterogeneous operating conditions and supports only a relative stress test, not high-accuracy external SOH or cross-laboratory generalization. Eighth, the method assumes that the first five charge records form a meaningful commissioning reference. Cells entering service after unobserved ageing or with corrupted initial records may violate that assumption.

Finally, PCHP always returns a bounded prediction and does not detect when abstention is preferable. A support or estimability gate could be added, but it would constitute a separate method claim. Such a gate should be specified without using the outcomes on which that new claim is evaluated and should be tested on data not used to design or select the gate; it does not inherently require newly collected laboratory data or an ever larger stack of datasets.

\section{Conclusion}

We introduced PCHP for cross-domain battery SOH estimation when target outcomes and future records are unavailable at prediction time. The method converts a protected source prediction into a one-sided causal state and projects a learned commissioning-referenced candidate into a recursively feasible harm tube. This yields prefix invariance, non-increasing output, $O(1)$ online state and computation, an exact per-record absolute-loss harm bound, and tight prefix-level $\ell_\infty$ nonexpansiveness of the complete causal cascade. The same analysis proves that the trajectory-constrained intersection is the exact online viability kernel, that a nontrivial unconditional zero-harm update is impossible, and that PCHP uniquely minimizes candidate distortion in that kernel. The loss-level analysis further identifies a reusable feasibility boundary: metric losses induce exact metric balls; bounded squared loss induces an exact baseline-dependent interval with the same recursive implementation; and squared loss over $\R$ admits no nontrivial finite outcome-uniform budget. Directionally weighted and time-varying corollaries provide asymmetric tubes, an exact lower-endpoint viability condition, and a tight minimum current budget. Strict nested evaluation across $12$ domains showed a domain-equal MAE reduction of $0.00486$; a matched candidate-free control and a protocol-locked retrospective budget path showed that the development gain was not confined to one constant budget-feasible shift or one budget value. A zero-additional-tuning audit retained the benefit across Ridge, HistGradientBoosting, and LightGBM, excluding an ExtraTrees-specific explanation within those audited families. A stronger unprotected comparator obtained substantially lower MAE but violated every domain budget, making the opportunity cost explicit. The outcome-blind, prospectively locked BaSyTec external validation improved all $45$ eligible cells and retained every harm check, while its fixed-shift diagnostic showed that this external dataset evaluates the contract rather than adaptive superiority. The NASA stress test exposed an incompatible absolute target regime. The contribution is therefore a deployable information and harm-control contract for learned battery-health updates, not a claim of accuracy dominance or universal SOH safety.

\section*{Declaration of generative AI and AI-assisted technologies in the manuscript preparation process}

During the preparation of this work, the author used OpenAI Codex to assist with research-code development and verification, literature-search organization, manuscript structuring, translation, and language editing. The author reviewed and independently verified all source claims, mathematical derivations, analyses, code outputs, figures, and text and takes full responsibility for the content of the work. All scientific figures were generated programmatically from the reported data and scripts; no generative image model was used to create or alter manuscript figures.

\section*{Data and code availability}

The exact development snapshot, BatteryLife\_Processed v11, is archived on Zenodo under CC BY~4.0 \citep{batterylife2026processed}. Its constituent domains originate from the cited providers and remain governed by their own access, attribution, and redistribution terms; consequently, no third-party raw archive is redistributed with this study. The $48$-cell Bayreuth degradation-path dataset is available from Zenodo under CC BY~4.0 \citep{roder2025degradationdata}; NASA PCoE Battery Data Set~5 is available from the official NASA Prognostics Data Repository \citep{saha2007nasa}. An executable review-lite reproducibility package accompanying the submission contains the complete transitive source-code closure, acquisition instructions, source and artifact hashes, the environment specification, pre-freeze records, aggregate per-domain and per-cell evidence, deterministic operator and statistical replay checks, and all figure-generation inputs and scripts. It reproduces the auditable evidence layer without redistributing third-party raw archives; complete source-model refitting and external replay require downloading provider data under their respective terms. The release licence for author-created code will be finalized with the institutional submission record.

\section*{Author information}

Author: Yuyang Wu. Affiliation, correspondence address, ORCID, funding, acknowledgements, and conflict-of-interest statements are intentionally left blank until the submission record is finalized.

\begingroup
\hbadness=2000
\bibliographystyle{plainnat}
\bibliography{references}
\endgroup

\end{document}
